\subsection{معرفی و انگیزه}
این مکانیزم بر مبنای مبانی نظری پیشرفته در ادبیات توجه و فشرده‌سازی \cite{eckart1936approximation} توسعه یافته است.

کوانتیزه‌سازی یکنواخت وزن‌های شبکه‌های عصبی عمیق به مقادیر با طول بیت پایین (\lr{Low-Bit Quantization})، موثرترین روش جهت کاهش ردپای حافظه و افزایش بازدهی مصرف انرژی است. با این وجود، در روش‌های سنتی کوانتیزه‌سازی با طول بیت ثابت، به تمام کانال‌ها و لایه‌ها نرخ بیت یکسانی اختصاص می‌یابد؛ در حالی که در ترنسفورمرهای بزرگ، توزیع انرژی و حساسیت کانال‌های وزنی به شدت ناهمگن است و بخش اندکی از ابعاد ویژگی سهم عمده‌ای در دقت پیش‌بینی مدل ایفا می‌کنند.

موتور کوانتیزه‌سازی پویا \lr{AstroQuantizer} با هدف غلبه بر این چالش طراحی گردیده است. این چارچوب نوآورانه از سیستم تخصیص پویای نرخ بیت بر اساس اطلاعات فیشر تجربی (\lr{Fisher-Adaptive Mixed-Precision Bit Allocation}) بهره می‌برد. در این مکانیزم، کانال‌هایی که حساسیت گرادیانی یا انرژی فعال‌سازی بالاتری دارند، طول بیت بیشتری (تا سقف ۸ بیت) دریافت کرده و کانال‌های کم‌اهمیت به بیت‌های پایین‌تر (تا کف ۲ بیت) نگاشته می‌شوند، در حالی که بودجه میانگین طول بیت کل لایه در سطح هدف دقیقا ارضا می‌گردد. همچنین با استفاده از ارزیاب گرادیان مستقیم (\lr{Straight-Through Estimator - STE})، جریان گرادیان در کل شبکه بدون انسداد حفظ می‌شود.
